\title{2021 年高考数学 上海春季卷}
\date{}
\maketitle

\section*{一、填空题}
本大题共12题，满分54分，第1～6题每题4分，第7～12题每题5分

\begin{question}[points=4]
已知等差数列 $\{a_n\}$ 的首项为3，公差为2，则 $a_{10}=$  .
\fillin[{$21$}]
\begin{solution}
\textbf{答案：}$21$

\textbf{分析：}由已知结合等差数列的通项公式即可直接求解．

\textbf{详解：}因为等差数列 $\{a_n\}$ 的首项为3，公差为2，则 $a_{10}=a_1+9d=3+9\times 2=21$．
\end{solution}
\end{question}

\begin{question}[points=4]
已知 $z=1-3i$，则 $|z-i|=$  .
\fillin[{$\sqrt{5}$}]
\begin{solution}
\textbf{答案：}$\sqrt{5}$

\textbf{分析：}由已知求得 $z-i$，再由复数模的计算公式求解．

\textbf{详解：}因为 $z=1-3i$，所以 $z-i=1+3i-i=1+2i$，则 $|z-i|=|1+2i|=\sqrt{1^2+2^2}=\sqrt{5}$．
\end{solution}
\end{question}

\begin{question}[points=4]
已知圆柱的底面半径为1，高为2，则圆柱的侧面积为  .
\fillin[{$4\pi$}]
\begin{solution}
\textbf{答案：}$4\pi$

\textbf{分析：}根据圆柱的侧面积公式计算即可．

\textbf{详解：}圆柱的底面半径为 $r=1$，高为 $h=2$，所以圆柱的侧面积为 $S_{\text{侧}}=2\pi rh=2\pi\times 1\times 2=4\pi$．
\end{solution}
\end{question}

\begin{question}[points=4]
不等式 $\frac{2x+5}{x-2}<1$ 的解集为  .
\fillin[{$(-7,2)$}]
\begin{solution}
\textbf{答案：}$(-7,2)$

\textbf{分析：}由已知进行转化 $\frac{x+7}{x-2}<0$，进行可求．

\textbf{详解：}由 $\frac{2x+5}{x-2}<1$ 得 $\frac{2x+5}{x-2}-1<0$，即 $\frac{x+7}{x-2}<0$，解得 $-7<x<2$．故答案为 $(-7,2)$．
\end{solution}
\end{question}

\begin{question}[points=4]
直线 $x=-2$ 与直线 $3x-y+1=0$ 的夹角为  .
\fillin[{$\frac{\pi}{6}$}]
\begin{solution}
\textbf{答案：}$\frac{\pi}{6}$

\textbf{分析：}先求出直线的斜率，可得它们的倾斜角，从而求出两条直线的夹角．

\textbf{详解：}直线 $x=-2$ 的斜率不存在，倾斜角为 $\frac{\pi}{2}$，直线 $3x-y+1=0$ 的斜率为 $\sqrt{3}$，倾斜角为 $\frac{\pi}{3}$，故直线 $x=-2$ 与直线 $3x-y+1=0$ 的夹角为 $\frac{\pi}{2}-\frac{\pi}{3}=\frac{\pi}{6}$．
\end{solution}
\end{question}

\begin{question}[points=4]
若方程组 $\begin{cases} a_1x+b_1y=c_1 \\ a_2x+b_2y=c_2 \end{cases}$ 无解，则 $\begin{vmatrix} a_1 & b_1 \\ a_2 & b_2 \end{vmatrix}=$  .
\fillin[{$0$}]
\begin{solution}
\textbf{答案：}$0$

\textbf{分析：}利用二元一次方程组的解的行列式表示进行分析即可得到答案．

\textbf{详解：}对于方程组 $\begin{cases} a_1x+b_1y=c_1 \\ a_2x+b_2y=c_2 \end{cases}$，有 $D=\begin{vmatrix} a_1 & b_1 \\ a_2 & b_2 \end{vmatrix}, D_x=\begin{vmatrix} c_1 & b_1 \\ c_2 & b_2 \end{vmatrix}, D_y=\begin{vmatrix} a_1 & c_1 \\ a_2 & c_2 \end{vmatrix}$，当 $D\neq 0$ 时，方程组的解为 $\begin{cases} x=\frac{D_x}{D} \\ y=\frac{D_y}{D} \end{cases}$，根据题意，方程组无解，所以 $D=0$，即 $\begin{vmatrix} a_1 & b_1 \\ a_2 & b_2 \end{vmatrix}=0$．
\end{solution}
\end{question}

\begin{question}[points=5]
已知 $(1+x)^n$ 的展开式中，唯有 $x^3$ 的系数最大，则 $(1+x)^n$ 的系数和为  .
\fillin[{$64$}]
\begin{solution}
\textbf{答案：}$64$

\textbf{分析：}由已知可得 $n=6$，令 $x=1$，即可求得系数和．

\textbf{详解：}由题意，$C_n^3>C_n^2$，且 $C_n^3>C_n^4$，所以 $n=6$，所以令 $x=1$，$(1+x)^6$ 的系数和为 $2^6=64$．
\end{solution}
\end{question}

\begin{question}[points=5]
已知函数 $f(x)=3^x+\frac{a}{3^x+1}$ $(a>0)$ 的最小值为5，则 $a=$  .
\fillin[{$9$}]
\begin{solution}
\textbf{答案：}$9$

\textbf{分析：}利用基本不等式求最值需要满足“一正、二定、三相等”，该题只需将函数解析式变形成 $f(x)=3^x+1+\frac{a}{3^x+1}-1$，然后利用基本不等式求解即可，注意等号成立的条件．

\textbf{详解：}$f(x)=3^x+\frac{a}{3^x+1}=3^x+1+\frac{a}{3^x+1}-1\geq 2\sqrt{a}-1=5$，所以 $a=9$，经检验，$3^x=2$ 时等号成立．
\end{solution}
\end{question}

\begin{question}[points=5]
在无穷等比数列 $\{a_n\}$ 中，$\lim\limits_{n\to\infty}(a_1-a_n)=4$，则 $a_2$ 的取值范围是  .
\fillin[{$(-4,0)\cup(0,4)$}]
\begin{solution}
\textbf{答案：}$(-4,0)\cup(0,4)$

\textbf{分析：}由无穷等比数列的概念可得公比 $q$ 的取值范围，再由极限的运算知 $a_1=4$，从而得解．

\textbf{详解：}因为无穷等比数列 $\{a_n\}$，所以公比 $q\in(-1,0)\cup(0,1)$，则 $\lim\limits_{n\to\infty}a_n=0$，所以 $\lim\limits_{n\to\infty}(a_1-a_n)=a_1=4$，所以 $a_2=a_1q=4q\in(-4,0)\cup(0,4)$．
\end{solution}
\end{question}

\begin{question}[points=5]
某人某天需要运动总时长大于等于60分钟，现有五项运动可以选择，如表所示，问有几种运动方式组合  .\begin{center}\begin{tabular}{|c|c|c|c|c|} \hline A运动 & B运动 & C运动 & D运动 & E运动 \\ \hline 7点-8点 & 8点-9点 & 9点-10点 & 10点-11点 & 11点-12点 \\ \hline 30分钟 & 20分钟 & 40分钟 & 30分钟 & 30分钟 \\ \hline \end{tabular}\end{center}
\fillin[{$23$种}]
\begin{solution}
\textbf{答案：}$23$种

\textbf{分析：}由题意知至少要选2种运动，并且选2种运动的情况中，$AB$、$DB$、$EB$的组合不符合题意，由此求出结果．

\textbf{详解：}由题意知，至少要选2种运动，并且选2种运动的情况中，$AB$、$DB$、$EB$的组合不符合题意；所以满足条件的运动组合方式为：$C_5^2+C_5^3+C_5^4+C_5^5-3=10+10+5+1-3=23$（种）．
\end{solution}
\end{question}

\begin{question}[points=5]
已知椭圆 $x^2+\frac{y^2}{b^2}=1(0<b<1)$ 的左、右焦点为 $F_1$、$F_2$，以 $O$ 为顶点，$F_2$ 为焦点作抛物线交椭圆于 $P$，且 $\angle PF_1F_2=45^\circ$，则抛物线的准线方程是  .
\fillin[{$x=1-\sqrt{2}$}]
\begin{solution}
\textbf{答案：}$x=1-\sqrt{2}$

\textbf{分析：}先设出椭圆的左右焦点坐标，进而可得抛物线的方程，设出直线 $PF_1$ 的方程并与抛物线联立，求出点 $P$ 的坐标，由此可得 $PF_2\perp F_1F_2$，进而可以求出 $PF_1$，$PF_2$ 的长度，再由椭圆的定义即可求解．

\textbf{详解：}设 $F_1(-c,0)$，$F_2(c,0)$，则抛物线 $y^2=4cx$，直线 $PF_1:y=x+c$，联立 $\begin{cases} y^2=4cx \\ y=x+c \end{cases}$，解得 $x=c$，$y=2c$，所以点 $P$ 的坐标为 $(c,2c)$，所以 $PF_2\perp F_1F_2$，又 $PF_2=F_2F_1=2c$，所以 $PF_1=2\sqrt{2}c$，所以 $PF_1+PF_2=(2+2\sqrt{2})c=2a=2$，则 $c=\sqrt{2}-1$，所以抛物线的准线方程为 $x=-c=1-\sqrt{2}$．
\end{solution}
\end{question}

\begin{question}[points=5]
已知 $\theta>0$，存在实数 $\varphi$，使得对任意 $n\in\mathbb{N}^*$，$\cos(n\theta+\varphi)<\frac{\sqrt{3}}{2}$，则 $\theta$ 的最小值是  .
\fillin[{$\frac{2\pi}{5}$}]
\begin{solution}
\textbf{答案：}$\frac{2\pi}{5}$

\textbf{分析：}在单位圆中分析可得 $\theta>\frac{\pi}{3}$，由 $\frac{2\pi}{\theta}\in\mathbb{N}^*$，即 $\theta=\frac{2\pi}{k}$，$k\in\mathbb{N}^*$，即可求得 $\theta$ 的最小值．

\textbf{详解：}在单位圆中分析，由题意可得 $n\theta+\varphi$ 的终边要落在图中阴影部分区域（其中 $\angle AOx=\angle BOx=\frac{\pi}{6}$），所以 $\theta>\angle AOB=\frac{\pi}{3}$，因为对任意 $n\in\mathbb{N}^*$ 都成立，所以 $\frac{2\pi}{\theta}\in\mathbb{N}^*$，即 $\theta=\frac{2\pi}{k}$，$k\in\mathbb{N}^*$，同时 $\theta>\frac{\pi}{3}$，所以 $\theta$ 的最小值为 $\frac{2\pi}{5}$．
\end{solution}
\end{question}

\section*{二、选择题}
本大题共4题，每题5分，共20分

\begin{question}[points=5]
下列函数中，在定义域内存在反函数的是 (\_\_\_\_\_\_\_\_\_\_)
\begin{choices}
  \item $f(x)=x^2$
  \item $f(x)=\sin x$
  \item $f(x)=2^x$
  \item $f(x)=1$
\end{choices}
\begin{solution}
\textbf{答案：}C

\textbf{分析：}根据函数的定义以及映射的定义即可判断选项是否正确．

\textbf{详解：}选项A：因为函数是二次函数，属于二对一的映射，根据函数的定义可得函数不存在反函数，A错误；选项B：因为函数是三角函数，有周期性和对称性，属于多对一的映射，根据函数的定义可得函数不存在反函数，B错误；选项C：因为函数是单调递增的指数函数，属于一一映射，所以函数存在反函数，C正确；选项D：因为函数是常数函数，属于多对一的映射，所以函数不存在反函数，D错误．故选C．
\end{solution}
\end{question}

\begin{question}[points=5]
已知集合 $A=\{x|x>-1,\,x\in\mathbb{R}\}$，$B=\{x|x^2-x-2\ge 0,\,x\in\mathbb{R}\}$，则下列关系中，正确的是 (\_\_\_\_\_\_\_\_\_\_)
\begin{choices}
  \item $A\subseteq B$
  \item $\\mathbb{R}\\setminus A\subseteq\\mathbb{R}\\setminus B$
  \item $A\cap B=\varnothing$
  \item $A\cup B=\mathbb{R}$
\end{choices}
\begin{solution}
\textbf{答案：}D

\textbf{分析：}根据集合的基本运算对每一选项判断即可．

\textbf{详解：}已知集合 $A=\{x|x>-1,x\in\mathbb{R}\}$，$B=\{x|x^2-x-2\ge 0,x\in\mathbb{R}\}$，解得 $B=\{x|x\ge 2\text{ 或 }x\le -1,x\in\mathbb{R}\}$，$\\mathbb{R}\\setminus A=\{x|x\le -1,x\in\mathbb{R}\}$，$\\mathbb{R}\\setminus B=\{x|-1<x<2\}$；则 $A\cup B=\mathbb{R}$，$A\cap B=\{x|x\ge 2\}$，故选D．
\end{solution}
\end{question}

\begin{question}[points=5]
已知函数 $y=f(x)$ 的定义域为 $\mathbb{R}$，下列是 $f(x)$ 无最大值的充分条件是 (\_\_\_\_\_\_\_\_\_\_)
\begin{choices}
  \item $f(x)$ 为偶函数且关于点 $(1,1)$ 对称
  \item $f(x)$ 为偶函数且关于直线 $x=1$ 对称
  \item $f(x)$ 为奇函数且关于点 $(1,1)$ 对称
  \item $f(x)$ 为奇函数且关于直线 $x=1$ 对称
\end{choices}
\begin{solution}
\textbf{答案：}C

\textbf{分析：}根据题意，依次判断选项：对于ABD，举出反例可得三个选项错误，对于C，利用反证法可得其正确．

\textbf{详解：}对于A，$f(x)=\cos\frac{\pi x}{2}+1$ 为偶函数且关于点 $(1,1)$ 对称，存在最大值，A错误；对于B，$f(x)=\cos(\pi x)$ 为偶函数且关于直线 $x=1$ 对称，存在最大值，B错误；对于C，假设 $f(x)$ 有最大值，设其最大值为 $M$，其最高点的坐标为 $(a,M)$，$f(x)$ 为奇函数，其图象关于原点对称，则 $f(x)$ 的图象存在最低点 $(-a,-M)$，又由 $f(x)$ 的图象关于点 $(1,1)$ 对称，则 $(-a,-M)$ 关于点 $(1,1)$ 对称的点为 $(2+a,2+M)$，与最大值为 $M$ 相矛盾，则此时 $f(x)$ 无最大值，C正确；对于D，$f(x)=\sin\frac{\pi x}{2}$ 为奇函数且关于直线 $x=1$ 对称，D错误．故选C．
\end{solution}
\end{question}

\begin{question}[points=5]
在 $\triangle ABC$ 中，$D$ 为 $BC$ 中点，$E$ 为 $AD$ 中点，则以下结论：①存在 $\triangle ABC$，使得 $\overrightarrow{AB}\cdot\overrightarrow{CE}=0$；②存在三角形 $\triangle ABC$，使得 $\overrightarrow{CE}//(\overrightarrow{CB}+\overrightarrow{CA})$；它们的成立情况是 (\_\_\_\_\_\_\_\_\_\_)
\begin{choices}
  \item ①成立，②成立
  \item ①成立，②不成立
  \item ①不成立，②成立
  \item ①不成立，②不成立
\end{choices}
\begin{solution}
\textbf{答案：}B

\textbf{分析：}设 $A(2x,2y)$，$B(-1,0)$，$C(1,0)$，$D(0,0)$，$E(x,y)$，由向量数量积的坐标运算即可判断①；$F$ 为 $AB$ 中点，可得 $(\overrightarrow{CB}+\overrightarrow{CA})=2\overrightarrow{CF}$，由 $D$ 为 $BC$ 中点，可得 $CF$ 与 $AD$ 的交点即为重心 $G$，从而可判断②．

\textbf{详解：}不妨设 $A(2x,2y)$，$B(-1,0)$，$C(1,0)$，$D(0,0)$，$E(x,y)$，① $\overrightarrow{AB}=(-1-2x,-2y)$，$\overrightarrow{CE}=(x-1,y)$，若 $\overrightarrow{AB}\cdot\overrightarrow{CE}=0$，则 $-(1+2x)(x-1)-2y^2=0$，即 $-(1+2x)(x-1)=2y^2$，满足条件的 $(x,y)$ 存在，例如 $(0,\frac{\sqrt{2}}{2})$，满足上式，所以①成立；② $F$ 为 $AB$ 中点，$(\overrightarrow{CB}+\overrightarrow{CA})=2\overrightarrow{CF}$，$CF$ 与 $AD$ 的交点即为重心 $G$，因为 $G$ 为 $AD$ 的三等分点，$E$ 为 $AD$ 中点，所以 $\overrightarrow{CE}$ 与 $\overrightarrow{CG}$ 不共线，即②不成立．故选B．
\end{solution}
\end{question}

\section*{三、解答题}
本大题共5题，共14+14+14+16+18=76分

\begin{question}[points=14]
四棱锥 $P-ABCD$，底面为正方形 $ABCD$，边长为4，$E$ 为 $AB$ 中点，$PE\perp$ 平面 $ABCD$．\begin{enumerate}\item[(1)] 若 $\triangle PAB$ 为等边三角形，求四棱锥 $P-ABCD$ 的体积；\item[(2)] 若 $CD$ 的中点为 $F$，$PF$ 与平面 $ABCD$ 所成角为 $45^\circ$，求 $PC$ 与 $AD$ 所成角的大小．\end{enumerate}
\begin{solution}
\textbf{答案：}$(1)\ \frac{32\sqrt{3}}{3}$；$(2)\ \arctan\frac{\sqrt{5}}{2}$

\textbf{分析：}（1）由 $V=\frac{1}{3}PE\cdot S_{\text{正方形}ABCD}$，代入相应数据，进行运算，即可；\newline（2）由 $PE\perp$ 平面 $ABCD$，知 $\angle PFE=45^\circ$，进而有 $PE=FE=4$，$PB=2\sqrt{5}$，由 $AD//BC$，知 $\angle PCB$ 或其补角即为所求，可证 $BC\perp$ 平面 $PAB$，从而有 $BC\perp PB$，最后在 $\text{Rt}\triangle PBC$ 中，由 $\tan\angle PCB=\frac{PB}{BC}$，得解．

\textbf{详解：}（1）因为 $\triangle PAB$ 为等边三角形，且 $E$ 为 $AB$ 中点，$AB=4$，所以 $PE=2\sqrt{3}$，又 $PE\perp$ 平面 $ABCD$，所以四棱锥 $P-ABCD$ 的体积 $V=\frac{1}{3}PE\cdot S_{\text{正方形}ABCD}=\frac{1}{3}\times 2\sqrt{3}\times 4^2=\frac{32\sqrt{3}}{3}$．\newline（2）因为 $PE\perp$ 平面 $ABCD$，所以 $\angle PFE$ 为 $PF$ 与平面 $ABCD$ 所成角为 $45^\circ$，即 $\angle PFE=45^\circ$，所以 $\triangle PEF$ 为等腰直角三角形，因为 $E$，$F$ 分别为 $AB$，$CD$ 的中点，所以 $PE=FE=4$，所以 $PB=\sqrt{PE^2+BE^2}=2\sqrt{5}$，因为 $AD//BC$，所以 $\angle PCB$ 或其补角即为 $PC$ 与 $AD$ 所成角，因为 $PE\perp$ 平面 $ABCD$，所以 $PE\perp BC$，又 $BC\perp AB$，$PE\cap AB=E$，$PE$、$AB\subset$ 平面 $PAB$，所以 $BC\perp$ 平面 $PAB$，所以 $BC\perp PB$，在 $\text{Rt}\triangle PBC$ 中，$\tan\angle PCB=\frac{PB}{BC}=\frac{2\sqrt{5}}{4}=\frac{\sqrt{5}}{2}$，故 $PC$ 与 $AD$ 所成角的大小为 $\arctan\frac{\sqrt{5}}{2}$．
\end{solution}
\end{question}

\begin{question}[points=14]
已知 $A$、$B$、$C$ 为 $\triangle ABC$ 的三个内角，$a$、$b$、$c$ 是其三条边，$a=2$，$\cos C=-\frac{1}{4}$．\begin{enumerate}\item[(1)] 若 $\sin A=2\sin B$，求 $b$、$c$；\item[(2)] 若 $\cos(A-\frac{\pi}{4})=\frac{4}{5}$，求 $c$．\end{enumerate}
\begin{solution}
\textbf{答案：}$(1)\ b=1,\ c=\sqrt{6}$；$(2)\ c=\frac{5\sqrt{30}}{2}$

\textbf{分析：}（1）由已知利用正弦定理即可求解 $b$ 的值；利用余弦定理即可求解 $c$ 的值．\newline（2）根据已知利用两角差的余弦公式，同角三角函数基本关系式可求得 $\cos A$，$\sin A$，$\sin C$ 的值，进而根据正弦定理可得 $c$ 的值．

\textbf{详解：}（1）因为 $\sin A=2\sin B$，可得 $a=2b$，又 $a=2$，可得 $b=1$，由余弦定理 $\cos C=\frac{a^2+b^2-c^2}{2ab}=\frac{2^2+1^2-c^2}{2\times 2\times 1}=-\frac{1}{4}$，解得 $c=\sqrt{6}$．\newline（2）因为 $\cos(A-\frac{\pi}{4})=\frac{\sqrt{2}}{2}(\cos A+\sin A)=\frac{4}{5}$，可得 $\cos A+\sin A=\frac{4\sqrt{2}}{5}$，又 $\cos^2 A+\sin^2 A=1$，解得 $\cos A=\frac{7\sqrt{2}}{10},\sin A=\frac{\sqrt{2}}{10}$ 或 $\sin A=\frac{7\sqrt{2}}{10},\cos A=\frac{\sqrt{2}}{10}$，因为 $\cos C=-\frac{1}{4}$，可得 $\sin C=\frac{\sqrt{15}}{4}$，$\tan C=-\sqrt{15}$，若 $\sin A=\frac{7\sqrt{2}}{10}$，$\cos A=\frac{\sqrt{2}}{10}$，可得 $\tan A=7$，则 $\tan B=-\tan(A+C)=\frac{\tan A+\tan C}{\tan A\tan C-1}=\frac{7-\sqrt{15}}{7\times(-\sqrt{15})-1}<0$，可得 $B$ 为钝角，这与 $C$ 为钝角矛盾，舍去，所以 $\sin A=\frac{\sqrt{2}}{10}$，由正弦定理 $\frac{a}{\sin A}=\frac{c}{\sin C}$，可得 $c=\frac{5\sqrt{30}}{2}$．
\end{solution}
\end{question}

\begin{question}[points=14]
（1）团队在 $O$ 点西侧、东侧20千米处设有 $A$、$B$ 两站点，测量距离发现一点 $P$ 满足 $|PA|-|PB|=20$ 千米，可知 $P$ 在 $A$、$B$ 为焦点的双曲线上，以 $O$ 点为原点，东侧为 $x$ 轴正半轴，北侧为 $y$ 轴正半轴，建立平面直角坐标系，$P$ 在北偏东 $60^\circ$ 处，求双曲线标准方程和 $P$ 点坐标．\newline（2）团队又在南侧、北侧15千米处设有 $C$、$D$ 两站点，测量距离发现 $|QA|-|QB|=30$ 千米，$|QC|-|QD|=10$ 千米，求 $|OQ|$（精确到1米）和 $Q$ 点位置（精确到1米，$1^\circ$）
\begin{solution}
\textbf{答案：}$(1)\ \frac{x^2}{100}-\frac{y^2}{300}=1,\ P(\frac{15\sqrt{2}}{2},\frac{5\sqrt{6}}{2})$；$(2)\ |OQ|\approx 19$ 米，$Q$ 点位置北偏东 $66^\circ$

\textbf{分析：}（1）求出 $a$，$c$，$b$ 的值即可求得双曲线方程，求出直线 $OP$ 的方程，与双曲线方程联立，即可求得 $P$ 点坐标；\newline（2）分别求出以 $A$、$B$ 为焦点，以 $C$、$D$ 为焦点的双曲线方程，联立即可求得点 $Q$ 的坐标，从而求得 $|OQ|$，及 $Q$ 点位置．

\textbf{详解：}（1）由题意可得 $a=10$，$c=20$，所以 $b^2=300$，所以双曲线的标准方程为 $\frac{x^2}{100}-\frac{y^2}{300}=1$，直线 $OP:y=\frac{\sqrt{3}}{3}x$，联立双曲线方程，可得 $x=\frac{15\sqrt{2}}{2}$，$y=\frac{5\sqrt{6}}{2}$，即点 $P$ 的坐标为 $(\frac{15\sqrt{2}}{2},\frac{5\sqrt{6}}{2})$．\newline（2）① $|QA|-|QB|=30$，则 $a=15$，$c=20$，所以 $b^2=175$，双曲线方程为 $\frac{x^2}{225}-\frac{y^2}{175}=1$；② $|QC|-|QD|=10$，则 $a=5$，$c=15$，所以 $b^2=200$，双曲线方程为 $\frac{y^2}{25}-\frac{x^2}{200}=1$，两双曲线方程联立，得 $Q(\frac{14400}{47},\frac{2975}{47})$，所以 $|OQ|\approx 19$ 米，$Q$ 点位置北偏东 $66^\circ$．
\end{solution}
\end{question}

\begin{question}[points=16]
已知函数 $f(x)=\sqrt{|x+a|-a}-x$．\begin{enumerate}\item[(1)] 若 $a=1$，求函数的定义域；\item[(2)] 若 $a\neq 0$，若 $f(ax)=a$ 有2个不同实数根，求 $a$ 的取值范围；\item[(3)] 是否存在实数 $a$，使得函数 $f(x)$ 在定义域内具有单调性？若存在，求出 $a$ 的取值范围．\end{enumerate}
\begin{solution}
\textbf{答案：}$(1)\ (-\infty,-2]\cup[0,+\infty)$；$(2)\ (0,\frac{1}{4})$；$(3)\ a\in(-\infty,-\frac{1}{4}]$

\textbf{分析：}（1）把 $a=1$ 代入函数解析式，由根式内部的代数式大于等于0求解绝对值的不等式得答案；\newline（2）$f(ax)=a\Leftrightarrow\sqrt{|ax+a|-a}=ax+a$，设 $ax+a=t\ge 0$，得 $a=t-t^2$，$t\ge 0$，求得等式右边关于 $t$ 的函数的值域可得 $a$ 的取值范围；\newline（3）分 $x\ge -a$ 与 $x<-a$ 两类变形，结合复合函数的单调性可得使得函数 $f(x)$ 在定义域内具有单调性的 $a$ 的范围．

\textbf{详解：}（1）当 $a=1$ 时，$f(x)=\sqrt{|x+1|-1}-x$，由 $|x+1|-1\ge 0$，得 $|x+1|\ge 1$，解得 $x\le -2$ 或 $x\ge 0$．所以函数的定义域为 $(-\infty,-2]\cup[0,+\infty)$．\newline（2）$f(ax)=\sqrt{|ax+a|-a}-ax$，$f(ax)=a\Leftrightarrow\sqrt{|ax+a|-a}=ax+a$，设 $ax+a=t\ge 0$，则 $\sqrt{t}-a=t$ 有两个不同实数根，整理得 $a=t-t^2$，$t\ge 0$，$a=-(t-\frac{1}{2})^2+\frac{1}{4}$，$t\ge 0$，当且仅当 $0<a<\frac{1}{4}$ 时，方程有2个不同实数根，又 $a\neq 0$，所以 $a$ 的取值范围是 $(0,\frac{1}{4})$．\newline（3）当 $x\ge -a$ 时，$f(x)=\sqrt{x+a-a}-x=\sqrt{x}-x=-(\sqrt{x}-\frac{1}{2})^2+\frac{1}{4}$，在 $[\frac{1}{4},+\infty)$ 上单调递减，此时需要满足 $-a\ge \frac{1}{4}$，即 $a\le -\frac{1}{4}$，函数 $f(x)$ 在 $[-a,+\infty)$ 上递减；当 $x<-a$ 时，$f(x)=\sqrt{-x-2a}-x$，在 $(-\infty,-2a]$ 上递减，因为 $a\le -\frac{1}{4}<0$，所以 $-2a>-a>0$，即当 $a\le -\frac{1}{4}$ 时，函数 $f(x)$ 在 $(-\infty,-a)$ 上递减．综上，当 $a\in(-\infty,-\frac{1}{4}]$ 时，函数 $f(x)$ 在定义域 $\mathbb{R}$ 上连续，且单调递减．
\end{solution}
\end{question}

\begin{question}[points=18]
已知数列 $\{a_n\}$ 满足 $a_n\ge 0$，对任意 $n\ge 2$，$a_n$ 和 $a_{n+1}$ 中存在一项使其为另一项与 $a_{n-1}$ 的等差中项．\begin{enumerate}\item[(1)] 已知 $a_1=5$，$a_2=3$，$a_4=2$，求 $a_3$ 的所有可能取值；\item[(2)] 已知 $a_1=a_4=a_7=0$，$a_2$、$a_5$、$a_8$ 为正数，求证：$a_2$、$a_5$、$a_8$ 成等比数列，并求出公比 $q$；\item[(3)] 已知数列中恰有3项为0，即 $a_r=a_s=a_t=0$，$2<r<s<t$，且 $a_1=1$，$a_2=2$，求 $a_{r+1}+a_{s+1}+a_{t+1}$ 的最大值．\end{enumerate}
\begin{solution}
\textbf{答案：}$(1)\ a_3=1$；$(2)\ q=\frac{1}{4}$；$(3)\ \frac{21}{64}$

\textbf{分析：}（1）根据 $a_n$ 和 $a_{n+1}$ 中存在一项使其为另一项与 $a_{n-1}$ 的等差中项建立等式，然后将 $a_1$，$a_2$，$a_4$ 的值代入即可；\newline（2）根据递推关系求出 $a_5$、$a_8$，然后根据等比数列的定义进行判定即可；\newline（3）分别求出 $a_{r+1}$，$a_{s+1}$，$a_{t+1}$ 的通项公式，从而可求出各自的最大值，从而可求出所求．

\textbf{详解：}（1）由题意，$2a_n=a_{n+1}+a_{n-1}$ 或 $2a_{n+1}=a_n+a_{n-1}$，所以 $2a_2=a_3+a_1$ 解得 $a_3=1$，$2a_3=a_2+a_1$ 解得 $a_3=4$，经检验，$a_3=1$．\newline（2）证明：因为 $a_1=a_4=a_7=0$，所以 $a_3=2a_2$ 或 $a_3=\frac{a_2}{2}$，经检验，$a_3=\frac{a_2}{2}$；所以 $a_5=\frac{a_3}{2}=\frac{a_2}{4}$ 或 $a_5=-a_1=-\frac{a_2}{2}$（舍），所以 $a_5=\frac{a_2}{4}$；$a_6=\frac{a_5}{2}=\frac{a_2}{8}$ 或 $a_6=-a_5=-\frac{a_2}{4}$（舍），所以 $a_6=\frac{a_2}{8}$；$a_8=\frac{a_6}{2}=\frac{a_2}{16}$ 或 $a_8=-a_6=-\frac{a_2}{8}$（舍），所以 $a_8=\frac{a_2}{16}$；综上，$a_2$、$a_5$、$a_8$ 成等比数列，公比为 $\frac{1}{4}$．\newline（3）由 $2a_n=a_{n+1}+a_{n-1}$ 或 $2a_{n+1}=a_n+a_{n-1}$，可知 $\frac{a_{n+2}-a_{n+1}}{a_{n+1}-a_n}=1$ 或 $\frac{a_{n+2}-a_{n+1}}{a_{n+1}-a_n}=-\frac{1}{2}$，由第（2）问可知，$a_r=0$，则 $a_{r-2}=2a_{r-1}$，即 $a_{r-1}-a_{r-2}=-a_{r-1}$，所以 $a_r=0$，则 $a_{r+1}=\frac{1}{2}a_{r-1}=-\frac{1}{2}(a_{r-1}-a_{r-2})=-\frac{1}{2}\cdot(-\frac{1}{2})^i\cdot 1^{r-3-i}\cdot(a_2-a_1)=-\frac{1}{2}\cdot(-\frac{1}{2})^i$，$i\in\mathbb{N}^*$，所以 $(a_{r+1})_{\max}=\frac{1}{4}$；同理，$a_{s+1}=-\frac{1}{2}\cdot(-\frac{1}{2})^j\cdot 1^{s-2-r-j}\cdot(a_{r+1}-a_r)=-\frac{1}{2}\cdot(-\frac{1}{2})^j\cdot\frac{1}{4}$，$j\in\mathbb{N}^*$，所以 $(a_{s+1})_{\max}=\frac{1}{16}$，同理，$(a_{t+1})_{\max}=\frac{1}{64}$，所以 $a_{r+1}+a_{s+1}+a_{t+1}$ 的最大值为 $\frac{21}{64}$．
\end{solution}
\end{question}

